% !TeX spellcheck = ru_RU
% !TEX root = vkr.tex

\section{Описание решения}
\label{sec:solution}

\subsection{Формат архива}

Разработанный архив начинается с заголовка фиксированного размера 2060~байт,
после которого располагается сжатый битовый поток. Заголовок позволяет
проверить тип файла, восстановить дерево и прекратить декодирование после
получения исходного количества байтов. Формат приведён в
таблице~\ref{tab:archive-format}.

\begin{table}[htbp]
    \centering
    \caption{Формат создаваемого архива.}
    \label{tab:archive-format}
    \begin{tabularx}{\textwidth}{@{}l l >{\raggedright\arraybackslash}X@{}}
        \toprule
        Поле & Размер & Назначение \\
        \midrule
        \texttt{magic} & 4 байта & сигнатура \texttt{HUFF} \\
        \texttt{size} & 8 байт & размер исходного файла, \texttt{uint64\_t}, little-endian \\
        \texttt{freq[256]} & 2048 байт & таблица частот возможных байтов \\
        \texttt{data} & переменный & битовый поток кодов Хаффмана \\
        \bottomrule
    \end{tabularx}
\end{table}

Таблица частот увеличивает размер малых архивов, но делает формат
самодостаточным: распаковка не требует отдельного словаря или внешних
метаданных. Представление размера файла и частот типом \texttt{uint64\_t}
исключает ограничение формата размером в 4~ГиБ.

\subsection{Архитектура реализации}

Реализация написана на языке~C стандарта C11 и разделена на следующие модули:
\begin{itemize}
    \item \texttt{main.c} разбирает аргументы командной строки и выбирает режим
          сжатия или распаковки;
    \item \texttt{huffman\_codec.c} подсчитывает частоты, формирует заголовок,
          строит таблицу кодов и выполняет декодирование;
    \item \texttt{huffman\_tree.c} строит и освобождает дерево Хаффмана,
          используя очередь с приоритетом на основе двоичной кучи;
    \item \texttt{bitio.c} реализует буферизованную запись и чтение отдельных
          битов в файловом потоке.
\end{itemize}

При сжатии файл читается сначала для подсчёта частот, затем повторно для записи
кодовых слов в битовый поток. При распаковке читается заголовок, по таблице
частот восстанавливается дерево и из потока выводится ровно число байтов,
заданное полем \texttt{size}. Специально обработаны пустой вход и файл,
состоящий из единственного различного символа: во втором случае символу
назначается однобитовый код.

\subsection{Инженерное оформление}

Сборка проекта описана средствами \CMake{}; при стандартной конфигурации для
исполняемого файла и тестов включается \AddressSanitizer{}. В репозитории
настроены процессы \GitHub{} Actions: сборка, модульные и интеграционные тесты,
контроль форматирования \texttt{clang-format}, а также статический анализ
\CodeQL{}. В README приведены команды сборки и использования, описание
формата архива и инструкция по запуску экспериментов. Исходный код
распространяется под лицензией MIT.
